\frontmatterpage
\unchapter{Anotācija}

Piemērotas programmatūras arhitektūras izvēle ir stratēģiski nozīmīgs inženiertehnisks un finansiāls lēmums maziem un vidējiem e-komercijas uzņēmumiem (MVU), kuriem jāsabalansē ātra funkcionalitātes izstrāde ar ierobežotu infrastruktūras budžetu un nelielu izstrādes komandu. Šajā bakalaura darbā veikts empīrisks salīdzinošs pētījums starp modulāro monolīto arhitektūru un sadalīto mikroservisu arhitektūru, izmantojot divas atvērtā pirmkoda references sistēmas: Microsoft \texttt{eShopOnWeb} un \texttt{eShopOnContainers}.

Darbs sastāv no piecām nodaļām. Pirmajā nodaļā aplūkoti programmatūras modularitātes, domēna vadītas projektēšanas un sadalīto sistēmu teorētiskie pamati. Otrajā nodaļā analizēti MVU e-komercijas platformu darbības ierobežojumi, mākoņdatošanas izmaksu struktūra un evolūcijas prasības. Trešajā nodaļā aprakstīts konteinerizētais eksperimentālais stends, slodzes testēšanas metodoloģija un pieci standartizēti funkcionālo izmaiņu scenāriji. Ceturtajā nodaļā apkopoti veiktspējas mērījumi un koda izmaiņu izplatības rezultāti. Piektajā nodaļā empīriskie dati sintezēti praktiskā arhitektūras izvēles ietvarā un izvērtēti pētījuma validitātes ierobežojumi.

Pētījuma metodoloģiskais pamats ietver zinātniskās literatūras analīzi, statisko pirmkoda atkarību un izmaiņu ietekmes novērtējumu piecos evolūcijas scenārijos, kontrolētus slodzes un resursu mērījumus ar Docker un \texttt{k6} rīkiem, kā arī salīdzinājumu ar publicētajiem industrijas migrācijas pētījumiem.

Empīriskie rezultāti parāda, ka modulārais monolīts patērē par 25,5\% mazāk operatīvās atmiņas miera stāvoklī (903,98 MiB pret 1\,213,60 MiB) un par 14,7\% mazāk atmiņas slodzes apstākļos (1\,151,35 MiB pret 1\,350,20 MiB), savukārt mikroservisu arhitektūra nodrošina 2,65 reizes zemāku 95.\ procentiles aizturi (12,30 ms pret 32,57 ms), pamatojot pakāpeniskas evolūcijas stratēģiju MVU uzņēmumiem.

\thesisscopeLV
